
\subsubsection{5.1 Existence: Trajectory Features Predict Minority Membership}

Per-sample loss trajectory features predict minority group membership with AUROC = 0.9452 $\pm$ 0.0072, exceeding the pre-specified threshold of 0.75 by a margin of 26\%.

\textbf{Table 1: Minority Group Prediction Performance}

\begin{table}[htbp]
\centering
\resizebox{\columnwidth}{!}{%
\begin{tabular}{ll}
\toprule
Metric & Value \\
\midrule
AUROC (mean $\pm$ std) & 0.9452 $\pm$ 0.0072 \\
Pre-specified threshold & 0.75 \\
Margin above threshold & +26.0\% \\
\bottomrule
\end{tabular}
}
\end{table}

This result indicates that trajectory features extracted from epochs 1-5 contain substantial information about minority group membership on this benchmark.

\subsubsection{5.2 Feature Analysis: Initial Loss Dominates}

Analysis of individual features reveals that initial loss ($L_{1})$ alone achieves AUROC = 0.9473, slightly exceeding the combined four-feature model.

\textbf{Table 2: Per-Feature AUROC}

\begin{table}[htbp]
\centering
\resizebox{\columnwidth}{!}{%
\begin{tabular}{ll}
\toprule
Feature & AUROC \\
\midrule
$L_1$ (Initial Loss) & 0.9473 \\
Slope & 0.8970 \\
Variance & 0.7242 \\
Convergence Time & 0.5259 \\
\bottomrule
\end{tabular}
}
\end{table}

The dominance of $L_1$ suggests that minority samples are distinguishable from the first training epoch. With pretrained ResNet-50, early layers already encode features that cause immediate conflict for minority samples where the spurious cue provides incorrect signal. This finding has practical implications: single-epoch screening may suffice for detection on this benchmark.

\subsubsection{5.3 Spurious-Specificity: Differential Attenuation Under GroupDRO}

The trajectory signal attenuates by 0.29 (31\% relative) under GroupDRO but only by 0.01 (1\% relative) under variance-matched random reweighting.

\textbf{Table 3: AUROC Across Training Regimes}

\begin{table}[htbp]
\centering
\resizebox{\columnwidth}{!}{%
\begin{tabular}{lll}
\toprule
Training Regime & AUROC (mean $\pm$ std) & $\Delta AUROC$ from ERM \\
\midrule
ERM (baseline) & 0.9436 $\pm$ 0.0123 & --- \\
GroupDRO & 0.6513 $\pm$ 0.0390 & -0.2923 \\
Random Reweighting & 0.9336 $\pm$ 0.0244 & -0.0100 \\
\bottomrule
\end{tabular}
}
\end{table}

GroupDRO specifically targets spurious correlation reliance by upweighting minority groups. The substantial attenuation under GroupDRO ($\Delta AUROC$ = 0.29, exceeding the 0.10 threshold) combined with minimal attenuation under random reweighting ($\Delta AUROC$ = 0.01, below the 0.05 threshold) is consistent with the hypothesis that trajectory divergence reflects spurious feature conflict rather than generic sample difficulty. The 29-fold difference in attenuation between the two interventions provides evidence for spurious-specificity.

Variance matching was verified: GroupDRO gradient variance was 0.015 and random reweighting gradient variance was 0.014, confirming that the control condition matched gradient variance without targeting spurious correlations.

\subsubsection{5.4 Negative Result: Curvature Timing Mechanism Not Supported}

The hypothesis that minority samples exhibit delayed curvature stabilization was not supported by the data.

\textbf{Table 4: Curvature Timing Analysis}

\begin{table}[htbp]
\centering
\resizebox{\columnwidth}{!}{%
\begin{tabular}{llll}
\toprule
Metric & Value & Threshold & Result \\
\midrule
Mean timing gap & 0.20 $\pm$ 0.40 epochs & $\geq$ 3 epochs & Not met \\
Seeds with gap $\geq$ 3 & 0/5 (0\%) & $\geq$ 70\% & Not met \\
\bottomrule
\end{tabular}
}
\end{table}

\textbf{Per-seed results:}

\begin{table}[htbp]
\centering
\resizebox{\columnwidth}{!}{%
\begin{tabular}{llll}
\toprule
Seed & Minority Median & Majority Median & Gap \\
\midrule
42 & 2.0 & 2.0 & 0.0 \\
123 & 2.0 & 2.0 & 0.0 \\
456 & 2.0 & 2.0 & 0.0 \\
789 & 3.0 & 2.0 & 1.0 \\
1011 & 2.0 & 2.0 & 0.0 \\
\bottomrule
\end{tabular}
}
\end{table}

Curvature stabilized at epochs 2-3 for both minority and majority samples across all seeds. The proposed timing gap of $\geq 3$ epochs was not observed in any seed. This negative result indicates that the discriminative signal comes from loss magnitude ($L_{1})$ rather than from temporal curvature dynamics. With pretrained models, convergence occurs quickly for all samples, leaving insufficient dynamic range to detect timing differences.

\subsubsection{5.5 Summary of Hypothesis Outcomes}

\begin{table*}[htbp]
\centering
\resizebox{\dimexpr 2\columnwidth+\columnsep\relax}{!}{%
\begin{tabular}{lllll}
\toprule
Hypothesis & Type & Threshold & Result & Outcome \\
\midrule
H-E1: Trajectory prediction & Existence & AUROC > 0.75 & 0.9452 & Supported \\
H-M1: Curvature timing & Mechanism & Gap $\geq$ 3 epochs, 70\% seeds & Gap = 0.20 epochs, 0\% & Not supported \\
H-M2: Spurious-specificity & Mechanism & $\Delta _GroupDRO$ > 0.10, $\Delta _Random$ < 0.05 & 0.29, 0.01 & Supported \\
\bottomrule
\end{tabular}
}
\end{table*}
